% TOP_PROBLEM: {"id": 2616, "title": "Kourovka Notebook Problem 21.107", "category_id": 17, "category": {"id": 17, "name": "group_theory", "display_name": "Group Theory", "description": "Problems about groups, group actions, representations, and related algebraic structures.", "slug": "group-theory", "order_index": 17, "created_at": "2026-07-31T15:26:25.671Z"}, "status": "open", "problem_number": "KOU-21.107", "set_id": 10, "statusQualification": "Makes the implicit nonempty-open-set convention explicit; retaining the literal empty open set would make expansiveness impossible."}
% FIRST_POSTED: 2026-10-01
% ENTRY_KIND: statement_only
% UnsolvedMath ID 2616; source catalogue code KOU-21.107
% !TeX program = lualatex
% Definition and sources only; this file is not an LLM solution attempt.
\documentclass[11pt,a4paper]{article}
\usepackage[margin=25mm]{geometry}
\usepackage{fontspec}
\setmainfont{lmroman10-regular.otf}[BoldFont=lmroman10-bold.otf,
 ItalicFont=lmroman10-italic.otf,BoldItalicFont=lmroman10-bolditalic.otf]
\usepackage{amsmath,amssymb,mathtools}
\usepackage{unicode-math}
\setmathfont{latinmodern-math.otf}
\AtBeginDocument{\let\setminus\smallsetminus}
\usepackage{xurl,hyperref}
\hypersetup{colorlinks=true,urlcolor=blue}
\setcounter{secnumdepth}{0}
\setlength{\parindent}{0pt}
\setlength{\parskip}{5pt}
\setlength{\emergencystretch}{3em}
\begin{document}
\section*{Kourovka Notebook Problem 21.107}\label{2616}
{\small UnsolvedMath ID 2616 · KOU-21.107 · Group Theory · Kourovka Notebook - New Problems, Issue 21}
\subsection{Definitions and mathematical statement}
A topological group is a group $G$ equipped with a topology for which multiplication $G\times G\to G$ and inversion $G\to G$ are continuous. A subset $D\subseteq G$ is dense if its closure is $G$, equivalently if it intersects every nonempty open subset of $G$.

Assume the underlying set of $G$ is countable and admits a partition into countably infinitely many dense subsets:
\[
 G=\bigcup_{j\ge1}D_j,\qquad D_i\cap D_j=\varnothing\ (i\ne j),
 \qquad\overline{D_j}=G\ (j\ge1).
\]
Is there a sequence $(F_n)_{n\ge1}$ of pairwise disjoint finite subsets of $G$ such that for every nonempty open set $U\subseteq G$,
\[
 \exists m=m(U)\quad\forall n>m:\quad F_n\cap U\ne\varnothing?
\]
Such a sequence is called expansive in the problem. The finite sets need not have equal sizes, and no bound on their sizes is requested.

The threshold may depend on the open set. Requiring intersection for every sufficiently large $n$ is stronger than requiring only that $\bigcup_nF_n$ be dense. The $F_n$ need not themselves partition $G$ and are not required to lie in the particular dense parts $D_j$. The countability assumption concerns the group itself, not a topological basis. The qualifier ``nonempty'' makes explicit the intended open-set convention, since no subset can intersect the empty open set.
\subsection{Short English statement}
Does every countable topological group partitionable into infinitely many dense subsets admit disjoint finite sets that eventually meet each nonempty open set?
\subsection{Sources}
\begin{itemize}
\item[S1] ULAM.AI and the UnsolvedMath Contributors, supplied problems.json, record 2616 (KOU-21.107), Kourovka Notebook - New Problems, Issue 21. Catalogue identity and source transcription; CC BY 4.0.
\url{https://huggingface.co/datasets/ulamai/UnsolvedMath}.
\item[S2] The Kourovka Notebook, 21st edition, new problems, Problem 21.107. Expanded formulation of the supplied catalogue statement.
\url{https://arxiv.org/pdf/1401.0300v47}.
\end{itemize}
\end{document}
